\defmodule {CONFIG}

Here, we define constants, elementary types, and configuration variables 
that are used throughout the package.

\bigskip\hrule
\code

DEFINITION MODULE CONFIG;

CONST
   MMaxDim = 16;
\endcode
  \tabb  Maximum number of dimensions of arrays and matrices 
    in module {\tt MULT}.
    Maximal order of a multiple recursive generator.
  \endtabb
\code

   BMaxDim = 48;
\endcode
 \tabb  Maximum number of dimensions for a lattice basis (module {\tt BASIS}).
 \endtabb
\code

   MaxCoord = 64;
\endcode
  \tabb  Maximum total number of coordinates that can be considered when
    constructing the generating vectors in module {\tt TESTLAT}.
  \endtabb
\code

   MaxDimProj = 8;
\endcode
  \tabb  Maximum number of dimensions for the projections over 
    non-successive indices, 
    that can be considered by the module {\tt TESTLAT}.
  \endtabb
\code

   MaxCat = 5;
\endcode
  \tabb    Maximum number of categories into which the
   set of dimensions can be partitioned.
  \endtabb
\code

   MaxNbGen = 256;
\endcode
 \tabb Maximum number of generators that can be stored in a list
  (in module {\tt SEARCH}). \endtabb
\code

   MaxJ = 5;
\endcode
 \tabb Maximum number of components $J$ for a combined generator. \endtabb
\code

TYPE
   MDim = [0..MMaxDim];
\endcode
 \tabb  A dimension or coordinate in module {\tt MULT}.
 \endtabb
\code

   BDim = [0..BMaxDim];
\endcode
 \tabb  A dimension or coordinate for a lattice basis.
 \endtabb
\code

   MCoord = [0..MaxCoord];
\endcode
 \tabb  A coordinate number that can be considered when
    constructing the generating vectors, and for the projections.
 \endtabb
\code

   DimProj = [0..MaxDimProj];
\endcode
 \tabb  Dimension of a projection.
 \endtabb
\code

   Categ = [0..MaxCat];
\endcode
  \tabb  A category of generators.
 \endtabb
\code

   Indexj = [1..MaxJ];
\endcode
 \tabb  Index of a component $j$ of a combined generator.  \endtabb
\code

   OutputType = (None, CurOut, CurOutTex);
\endcode
 \tabb  Kind of output to be produced.
  {\tt None} means that nothing is written to the current output.
  {\tt CurOut} and {\tt CurOutTex} mean that the results will be
  written to the current default output. The difference is that for the
  latter, they are written in \LaTeX\ format.
 \endtabb
\code

END CONFIG.
\endcode
